%% Two facing pages, outer ornament, class-independent headings and a reserved callout.
\documentclass[11pt,a4paper,twoside]{article}
\renewcommand\rmdefault{cmr}\renewcommand\sfdefault{cmss}\renewcommand\ttdefault{cmtt}
\usepackage[T1]{fontenc}
\usepackage[inner=22mm,outer=68mm,bindingoffset=3mm,top=32mm,bottom=30mm,headsep=8mm]{geometry}
\usepackage[grown,sections,novisitors,side=outer,style=chapter]{florilegium}
\pagestyle{plain}
\newcommand\sampletext{The garden follows the outer edge of the book. Page parity selects
the side, while the actual text block includes the binding allowance and any page offsets.
The text remains upright; only the ornament is reflected.\par
An editor can choose a quiet margin with one main focal flower, or a chapter treatment
with two larger focal motifs. Density, focal scale and asymmetry can then be adjusted
individually. These settings describe a composition, not a measured percentage of ink.\par
Reserved content is considered while planning branches and choosing motifs. A final
clipping boundary keeps the reserved area clear even when no route around it is possible.
The protected content itself is typeset normally.}
\begin{document}
\section{A chapter treatment on the right}
\sampletext
\begin{florilegiumprotect}[3mm]
  \medskip\hrule\medskip
  This unbreakable protected box can also be placed inside a figure or table float.
  Its reservation travels with the content when the float moves.
  \medskip\hrule\medskip
\end{florilegiumprotect}
\clearpage
\florilegiumsetup{style=quiet}
\section{A quiet margin on the left}
\sampletext
\florilegiumreserve[3mm]{12mm}{112mm}{30mm}{32mm}
\begin{tikzpicture}[remember picture,overlay]
  \begin{scope}[shift={(current page.south west)}]
    \filldraw[fill=black!4,draw=black!25] (1.2,11.2) rectangle (4.2,14.4);
    \node[align=left,text width=24mm,font=\small] at (2.7,12.8)
      {An intentional callout. The plant grows around this reserved area.};
  \end{scope}
\end{tikzpicture}
\end{document}
